\documentclass[a4paper, 11pt]{article}
\usepackage{fullpage}
\usepackage{mathtools}
\usepackage{amssymb}

\title{
  {Applied Analysis Seminar Notes}
  }
\begin{document}
\maketitle
\section{Navier-Stokes Equations}
Navier-Stokes equations are non-linear PDEs that govern the flow of fluids
and gases. There are many different types of flow regimes and depending on the modelling
needs one can make different simplification assumtions to help with the computation.

In the most general case there are three conservation laws at play. The conservation of
momentum, conservation of mass and conservation of energy. The energy conservation law is
important for cases where heating effects play a role.

The goal of this seminar is to gain fundamental understanding of these equations and solution
approaches. Since the heating effects add further complications to an already complex system,
only cases without these effects are considered here.

In this seminar project the benchmark model of a lid-driven cavity in 2D is considered.
It may appear as a simple model, yet it hides plethora of depth to gain a solid understanding
of the challenges one can face in simulating intricate domains and different flow regimes.

The Navier-Stokes equations is a coupled system of PDEs  of the following form:
\begin{equation}
\begin{cases}
  \frac{\partial u}{\partial t} - \operatorname{div}\left[\nu\left(\nabla u + \nabla u^T\right)\right] + (u \cdot \nabla)u + \nabla p = f, & x \in \Omega, \, t > 0, \\[5pt]
  \operatorname{div} u = 0, & x \in \Omega, \, t > 0,
\end{cases}
\end{equation}
This is the general form where $\Omega$ is the calculation domain. In this case, it is a unit square
domain of the form $\Omega = (0, 1)\times(0, 1)$. Now, there are no external foces acting on the
fluid inside the domain hence the right hand side of the first equaiton is thus zero.
\begin{equation}
\begin{cases}
  \frac{\partial u}{\partial t} - \operatorname{div}\left[\nu\left(\nabla u + \nabla u^T\right)\right] + (u \cdot \nabla)u + \nabla p = 0, & x \in \Omega, \, t > 0, \\[5pt]
  \operatorname{div} u = 0, & x \in \Omega, \, t > 0,
\end{cases}
\end{equation}
By non-dimensionalising the equaitons one gets the form:
\begin{equation}
\begin{cases}
  \frac{\partial u}{\partial t} - \frac{1}{\mathrm{Re}}\operatorname{div}\left[\left(\nabla u + \nabla u^T\right)\right] + (u \cdot \nabla)u + \nabla p = 0, & x \in \Omega, \, t > 0, \\[5pt]
  \operatorname{div} u = 0, & x \in \Omega, \, t > 0,
\end{cases}
\end{equation}
where Re is the Reynold's number $\mathrm{Re} = \frac{UL}{\nu}$. This is the only meaningfull
parameter in this case. Depending on the size the number turbulences may or may not appear.
\end{document}
